\chapter{问题、对象与统一结论}

本书研究一个明确的算法结构：在第 $k$ 个 outer step 用当前策略 $\pi_k$ 采样或校正访问测度，构造教师相对当前策略的信号，然后在固定 anchor 上进行若干 inner 更新，最后把 student 同步为 $\pi_{k+1}$。研究重点是同步顺序、镜像步长 $\eta$、Natural-GD 近似以及 EMA anchor 对动力学的改变。

\section{三层对象}

给定 prompt 或状态 $s$，记动作空间为有限集合 $\mathcal A$，学生策略为 $\pi$, teacher 为 $T$，行为策略或 outer anchor 为 $\mu$。需要区分三个层次：

\begin{enumerate}
    \item \textbf{分布层：}在 simplex 上求解 KL 镜像子问题，得到精确的目标分布；
    \item \textbf{参数层：}用神经网络和有限 inner steps 逼近该目标，产生投影误差；
    \item \textbf{访问测度层：}rollout 只覆盖实际到达的 prefix，决定分布层结论能否汇总为 sequence-level 结论。
\end{enumerate}

若把三层混为一层，就会错误地把“inner loss 已收敛”当成“全局策略已经等于 teacher”，也会把固定的 $\eta$ 误解成固定的 rollout 数据。正确的 outer pipeline 是
\[
\pi_k\;\longrightarrow\;\text{rollout 或 replay 校正}\;\longrightarrow\;
\text{冻结 }(\pi_k,T,r_k)\;\longrightarrow\;\text{inner solve}\;\longrightarrow\;
\pi_{k+1}\;\longrightarrow\;\text{重新采样}.
\]

\section{精确分布动力学的起点}

在固定 prefix、共同 support 和精确分布优化的条件下，OPD 的镜像目标为
\[
 p_{k+1}^{\star}(a\mid s)=
 \frac{\pi_k(a\mid s)^{1-\eta}T(a\mid s)^\eta}
 {\sum_b\pi_k(b\mid s)^{1-\eta}T(b\mid s)^\eta}.
\]
因此 $\eta=1$ 是一次 teacher 投影，$0<\eta<1$ 是 teacher 与 anchor 的几何插值。令 $\rho_k=\log(\pi_k/T)$，则去掉归一化常数后
\[
\rho_{k+1}(a\mid s)-\rho_{k+1}(a'\mid s)
=(1-\eta)\bigl(\rho_k(a\mid s)-\rho_k(a'\mid s)\bigr).
\]
这给出精确的 projective contraction，而不是参数空间中未经说明的普通线性插值。

有限 inner steps 时，实际迭代多出一个投影误差。若该误差在 projective 范数下为 $\varepsilon_k$，则
\[
 \Delta_{k+1}\le (1-\eta)\Delta_k+\varepsilon_k,
 \qquad
 \limsup_k\Delta_k\le \frac{\sup_k\varepsilon_k}{\eta}
\]
（常数 $\eta$ 且误差有界）。故减小 $\eta$ 会增加 outer 收敛步数，也会放大固定 inner 误差的稳态影响。

\section{从局部结论到完整算法的条件}

上述结论需要共同 support、有限 log-ratio 和正确的 prefix 测度。vanilla OPD 若只用一批旧 rollout、丢弃 anchor 依赖、又不控制 teacher support，则没有无条件的全局收敛、单调真实回报或有限方差结论。可证明的路线是先证明局部 Gibbs 与 projective 结果，再加入 inner 误差、occupancy mismatch、teacher 误差和采样噪声。

\section{符号、时钟与比较对象}
\begin{table}[htbp]
\centering
\begin{tabular}{p{0.21\textwidth}p{0.71\textwidth}}
\toprule 符号 & 含义\\\midrule
$s,a,\mathcal A,H$ & prefix、token、有限动作集合、response 的有限最大长度。\\
$T,p_\theta,b,\bar\pi$ & teacher、actor、生成数据的 behavior、镜像 anchor。\\
$k,m,j$ & outer iteration、该 outer 内的 inner update、EMA 更新事件。\\
$\eta,\tau=\eta^{-1}$ & 镜像步长与其逆温度；学习率另用 $\alpha$ 或 $\gamma$。\\
$\nu,d_p,d_b$ & 固定训练 prefix 测度、actor 和 behavior 的访问测度。\\
$G_\eta(b),\varepsilon_k$ & 精确镜像目标、实际 inner 解对该目标的误差。\\
$F,\epsilon_{\mathrm{NG}}$ & 参数 Fisher、inner 实际移动的 KL 半径。\\
$\beta,h,T_{\mathrm{ema}}$ & EMA 正则强度、一次同步保留旧 anchor 的系数、同步周期；第 5 章明确与论文原记号的转换。\\
\bottomrule
\end{tabular}
\end{table}

固定 $\eta$ 是每次 outer 都采用同一个近端强度。
rollout 是否更新则由采样时钟决定；$\eta$ 的数值没有定义采样时钟。
在一次 outer 内把固定子问题求至收敛，得到 $G_\eta(\pi_k)$；
hard sync 后下一次目标为 $G_\eta(\pi_{k+1})$。
因此训练不是在 $\pi_0$ 的同一批数据上永远重复同一个子问题。
MDPO 的 on-policy 与 off-policy 实现还有不同的同步频率，第二章逐项说明。

\section{两动作例子贯穿整个推导}
令 $b=(1/2,1/2)$、$T=(4/5,1/5)$。
teacher 的 odds 为 $4$，行为 odds 为 $1$，故镜像目标的 odds 是 $4^\eta$。
当 $\eta=1/2$ 时，第一次精确 outer 解为 $(2/3,1/3)$，第二次精确解的 odds 为 $4^{3/4}$。
经过 $k$ 次精确 hard sync，odds 为
\[
 \frac{p_k(1)}{p_k(2)}=4^{1-(1-\eta)^k}.
\]
inner 未收敛时，设第 $k$ 次 log-odds 求解误差为 $e_k$，则
\[
 z_{k+1}=(1-\eta)z_k+e_k,
 \qquad z_k=\log\frac{p_k(1)}{p_k(2)}-\log4.
\]
这条标量式解释首轮误差怎样被后续 outer 衰减，也解释反复产生同向求解误差为何形成地板。
第 4 章将它推广到多动作、变步长和有限 inner 次数。

$\eta=1$ 时精确目标一次等于 $T$，但参数训练只把 actor 推向这个目标。
未访问 prefix、不充分 inner 求解或不可表示的 teacher 都会留下误差。
这里的“到 teacher”是分布优化的解，不能把它等同于一次参数梯度更新。

\section{正文的论证依赖}
第二章从 Bregman 三点恒等式推到 PMD、MDPO 与 Natural-GD。
第三章用 Gibbs 变分式定义 MirrorOPD，并证明中心化残差与局部目标的关系，然后区分 sequence 与 prefix 测度。
第四章研究 hard sync、finite inner 和 $\eta$，先给精确动力学，再逐项加入误差。
第五章保留 EMA 原论文的参数和时钟，分别分析 anchor 滞后与 actor 收敛。
第六章完整展开 TOP-D 的定理证明并比较 vanilla OPD 和 PMD-mean。
第七章从概率距离与测度误差导出方差、性能下界及其附加条件。
第八章将这些命题转成可检验算法和实验；第九章说明与前序工作的关系。

所有“全局”结论都注明其量化范围：固定 prefix、全部有限 prefix 树，或神经网络验证集合。
对固定 teacher 的蒸馏收敛与对真实环境奖励的策略改进分别证明。
这样每个结论可以沿着明确的假设追溯到一个实际训练对象。
